\documentclass[11pt]{article}
\usepackage[a4paper,margin=25mm,marginparwidth=18mm]{geometry}
\usepackage{amsmath,amssymb,amsthm,booktabs,longtable,array,hyperref}
\setlength{\emergencystretch}{2em}
\newtheorem{theorem}{Theorem}\newtheorem{proposition}[theorem]{Proposition}
\newtheorem{lemma}[theorem]{Lemma}\newtheorem{corollary}[theorem]{Corollary}
\newcommand{\gen}[1]{\langle#1\rangle}
\newcommand{\Sym}{\operatorname{Sym}}\newcommand{\Alt}{\operatorname{Alt}}
\newcommand{\PSL}{\operatorname{PSL}}\newcommand{\PGL}{\operatorname{PGL}}
\newcommand{\PSU}{\operatorname{PSU}}\newcommand{\Sz}{\operatorname{Sz}}
\newcommand{\Sp}{\operatorname{Sp}}\newcommand{\Aut}{\operatorname{Aut}}
\newcommand{\note}[1]{\marginpar{\raggedright\scriptsize #1}}
\title{Rank spectra, permutation orbits and almost simple groups}
\author{An attributed synthesis of the fixed classification survey}
\date{14 October 2026}
\begin{document}\maketitle

We consider the twenty-one-page edition of Leemans's survey
\cite{survey}, arXiv:1910.08843v2, 11 May 2020. It contains thirty-one
numbered theorems, one corollary, six tables, one diagram figure,
fifteen problems and three conjectures. A survey has a different
logical contract from a classification paper: most of its theorems
are statements of earlier results. We therefore distinguish the proofs
acquired in this project from the published classification theorems
that remain explicit inputs. This distinction applies to every row
of the accompanying clause register. Recovering this synthesis does
not confer recovery credit on any other original paper.

Our companion \emph{Three elementary interfaces for a survey of string C-groups} proves the full rank-three intersection theorem,
the symmetric rank constructions with their exact corrected types,
and two counterexamples to the literal survey Theorem~5.7.
Here we assemble those results with the acquired uniform theories,
record the exact numerical conventions, and give an explicit disposition
for every source claim and question. We use ``acquired'' only for a complete proof
with its recorded independent review and acceptance; it does not mean
an unproved statement inferred from a matching numerical answer.

\section{Three principles and their actual hypotheses}
For an ordered involution marking $\rho=(\rho_0,\ldots,\rho_{r-1})$
write $G_I=\gen{\rho_i:i\in I}$. Its intersection property is
\[
                  G_I\cap G_J=G_{I\cap J}\qquad(I,J\subseteq\{0,\ldots,r-1\}).
\]
The string condition requires distant marks to commute. In all uses
of a subgroup, the marking is the actual subgroup marking. Passing
to an orbit or a block action is a homomorphism and does not, by
itself, preserve the intersection property.

The following three acquired principles account for much of the survey.
First, a faithful action of a finite group $H$ is the disjoint union
of coset actions $H/K$, with conjugacy classes of stabilizers and their
multiplicities as complete invariants. Its kernel is the intersection
of the corresponding cores. Stabilizers of a partial marking act on
possible next marks, so successive orbit selection gives a complete
census when every admissible extension is tested. The acquired finite
census proofs establish this induction, the full intersection test,
and the bound $r\le\lfloor\log_2|G|\rfloor$ used for termination.

Second, the acquired restriction-kernel theorem controls adjoining a
leaf to a permutation marking. It tests the actual kernels of prefix
restrictions, rather than assuming an extension preserves intersections.
For high-rank symmetric markings, the pure-cut theorem strengthens
this to an invertible subdivision operation, compatible with isomorphism
and reversal. It proves stability, not merely the construction of
additional examples.

Third, the acquired contraction theorem says that an irreducible
string C-marking of rank $r\ge4$ contracts to
\[
              (\rho_1,\rho_0\rho_2,\rho_3,\ldots,\rho_{r-1})
\]
whenever $\rho_0\in\gen{\rho_0\rho_2,\rho_3}$. If
$p_j=|\rho_{j-1}\rho_j|$, oddness of $p_3$ suffices: the corresponding
odd power of $(\rho_0\rho_2)\rho_3$ recovers $\rho_0$.
The complete acquired proof pays generation and every intersection.
Repeatedly testing $p_3,p_4,\ldots,p_{t+2}$ gives all ranks from
$r-t$ through $r$, for $0\le t\le r-3$.
\note{Reuse}
No step tests $p_2$. The extra initial test in survey Corollary~1 can
leave its displayed maximum undefined. The corrected iteration permits
zero steps without an oddness assumption and implies every defined
instance of the printed conclusion.

\begin{lemma}\label{connected}
A C-marking of an almost simple group has connected Coxeter diagram.
\end{lemma}
\begin{proof}
Let $S$ be its nonabelian simple socle. A disconnected diagram and
the intersection property would give a direct product $G=A\times B$
of nontrivial normal subgroups. Since $A\cap S$ is normal in $S$, it
is either $S$ or $1$. In the first case $B$ centralizes $S$; in the
second $[A,S]\le A\cap S=1$, so $A$ centralizes $S$. But
$S\le G\le\Aut(S)$ has $C_G(S)=1$. Both alternatives are impossible.
\end{proof}

For rank three, the companion proves the theorem of Conder and
Oliveros: a finite nontrivial group with no nontrivial cyclic normal
subgroup, generated by involutions $a,b,c$ with $ac=ca$, has the full
intersection property. The argument includes the exceptional
Klein-four intersection of its two dihedral subgroups. Thus the
rank-three positive simple-group classification reduces to involution
generation, subject to its genuinely substantial published input.
We retain the Nuzhin--Mazurov classification, with the Macaj--Jones
corrections, under its original attribution. The acquired simple-group
obstruction separately proves that the twelve exceptional families
listed in survey Theorem~3.1 have no string C-marking of any rank.

\section{Symmetric and alternating groups}
The adjacent transpositions along a path generate $S_n$. A selected
set generates the direct product of the symmetric groups on the path
components; intersections correspond to common selected edges. Thus
this is a C-marking. Contraction gives every rank $3\le r\le n-1$.
The companion computes the exact types:
\[
 \begin{cases}
 [3^{n-2}],&r=n-1,\\
 [n-r+2,6,3^{r-3}],&3\le r\le n-2.
 \end{cases}                                                    \tag{1}
\]
\note{Correction}
The second line cannot include the first case, as the literal survey
Theorem~4.3 does: already $S_4$ has no element of order six.
This corrects the type, not the rank-existence theorem.

The acquired symmetric stability theorem proves that the number of
markings of rank $n-k$, modulo isomorphism and duality, is independent
of $n$ for $n\ge2k+3$. Its first six values are
\[
                           1,\ 1,\ 7,\ 9,\ 35,\ 48.           \tag{2}
\]
The proof is the complete pure-cut bijection, together with the
complete rooted-wing calculation at the boundary. It resolves survey
Conjecture~1 as a consequence of the later original \cite{stability};
this is a recovered previous result, not a new solution to an open
conjecture. The separate acquired diagram classification for
$1\le k\le4$ matches all eighteen diagrams of the survey's Figure~1.
The original TeX arrows, parallel edges and ellipses were parsed with
the unchanged accepted source adapter. Exact point conjugacies to all
eighteen reviewed family models are retained. The uniform subdivision
proof then covers every parameter, including the smaller boundary
where the printed long ellipses do not literally fit.

For $A_n$, the acquired complete maximal-rank theorem gives
\[
 \max r=\begin{cases}
 3,&n=5,\\4,&n=9,\\5,&n=10,\\6,&n=11,\\
 \lfloor(n-1)/2\rfloor,&n\ge12.
 \end{cases}                                                   \tag{3}
\]
There is no C-marking for $n=3,4,6,7,8$. The acquired construction
and contraction proof gives every rank $3\le r\le\lfloor(n-1)/2\rfloor$
when $n\ge12$. In particular it implies the existence conclusion
``some alternating group for every rank'' in survey Theorem~4.6.
Its two more precise displayed families,
\[
 A_{2r+1}:[10,3^{r-2}]\ (r\ge4\text{ even}),\qquad
 A_{2r+3}:[10,3^{r-4},6,4]\ (r\ge5\text{ odd}),                \tag{4}
\]
remain explicitly the constructions of Fernandes--Leemans--Mixer
\cite{alt2012}. Their exact types and diagrams are not inferred from
an existence proof for another marking. The first is already an input
of the acquired alternating interval proof; its other maximal seed
is different from the second family in (4). This distinction prevents
an accidental transfer of an unproved Schlaefli type.

\subsection*{Pairs of commuting involutions}
\note{Reuse}
Two distinct nonidentity commuting involutions generate $V=C_2^2$.
Let $f,x,y,z,t$ be the multiplicities of its orbits with stabilizers
$V,\gen a,\gen c,\gen{ac},1$, respectively. A degree-$N$ pair is
faithful precisely when
\[
 N=f+2(x+y+z)+4t,\qquad t>0\text{ or at least two of }x,y,z>0.    \tag{5}
\]
Exchanging the involutions exchanges $x,y$. Both involutions are even
precisely when $x\equiv y\equiv z\pmod2$. Thus counting (5) with
$x\le y$ proves the two formulas of survey Theorems~4.4 and~4.10,
by the complete acquired multiplicity calculation. At degree six,
full automorphism fusion replaces the permutation-conjugacy answers
by five classes for $S_6$ and one for $A_6$. We retain that complete
finite automorphism calculation, not an exception inferred from the
source table. All fourteen rows of survey Table~2 agree with these
formulas; a separate literal orbit count agrees for every degree at
most fifty. The uniform proof is the orbit classification (5).

\section{Linear and classical groups}
The complete acquired subgroup-separator proofs recover the following
string classifications: $\PSL_2(q)$ has rank at most four, with rank
four precisely at $q=11,19$, yielding respectively the unique
11-cell and 57-cell; $\PGL_2(q)$ has rank at most four, with rank
four precisely at $q=5$, yielding the simplex. The standard geometry
and presentation identifications remain named inputs of those proofs.
The rank-three existence classification of Sjerve--Cherkassoff is
retained: $\PSL_2(q)$ occurs exactly for $q\notin\{2,3,7,9\}$,
and $\PGL_2(q)$ occurs exactly for $q\ne2$.

The larger classification in survey Theorem~5.6 is a separate input:
Connor--De Saedeleer--Leemans \cite{almost2} give the group
classification and rank assertions for almost simple groups with $\PSL_2$ socle.
Its small cases, even-field cases, odd fields of odd exponent, and
Baer extensions are retained exactly in the clause register. Neither
the string results for $\PSL_2,\PGL_2$ alone nor the project's finite
semilinear chiral census proves that infinite theorem. We retain its
published proof cost rather than hiding it in a label such as
``standard classification''.

\note{Correction}
The unnamed hypermap assertion at the start of Section~5.4 is false.
The separately checked companion gives explicit matrices in
$\operatorname{PSL}_2(11)$ defining a regular hypermap of type $(6,6,5)$.
Its three distinguished involutions generate the full group, while
$|\gen{a,b}\cap\gen{b,c}|=4>2=|\gen b|$.
This example lies in the hyperbolic, inner-reflexible range actually
classified by Conder--Poto\v cnik--\v Sir\'a\v n~\cite{hypermaps}.
The missing condition in the survey's appeal to its Theorem~3.2 is the
commutation of two distinguished involutions. The valid string theorem
and the original hypermap classification are unaffected. This is a
correction of an unnamed inference, not a newly settled open question.

\note{Correction}
Survey Theorem~5.7, stated for an arbitrary subgroup of
$\operatorname{GL}_3(q)$, is false. The companion's faithful $S_4$
action on $\{x\in\mathbb F_q^4:\sum x_i=0\}$ has the simplex
C-marking in characteristic two, contradicting the unqualified claim
that $q$ is odd. Over $\mathbb F_5$ the same action is absolutely
irreducible, but its order $24$ is smaller than
$|\Omega_3(5)|=60$. It therefore also refutes the displayed lower
containment, even after adding irreducibility. The body Theorem~3.3
of the cited Brooksbank--Vicinsky manuscript states form preservation,
without that lower containment.

The acquired corrected projective theorem suffices for the intended
application. An irreducible projective string C-group in dimension
three has rank at most four, is in odd characteristic, and preserves
a nondegenerate form after lifting. The natural projective actions
of $\PSL_3(q)$ and $\PGL_3(q)$ are irreducible. In odd characteristic
\[
       \frac{|O_3(q)|}{|\PSL_3(q)|}
       =\frac{2\gcd(3,q-1)}{q^2(q^3-1)}<1,
\]
so neither full projective group can occur. This is an actual form
obstruction, not a lower-containment assertion about every subgroup.
For $\PSL_4(q)$, survey Theorem~5.8 retains the published odd-field
rank-four construction of Brooksbank--Leemans \cite{linear4}; the
acquired simple-group obstruction excludes every even field.
The previously recorded project work on the maximum rank is kept
separate in Section~\ref{questions}.

For $q=2^k\ge4$, the acquired quadratic-path construction gives both
$O^+(2m,q)$ and $O^-(2m,q)$ in rank $2m$, and $\Sp(2m,q)$ in
rank $2m+1$, including $m=1$. Its generation and intersection proofs
use supported symmetry groups on actual subspaces. Free nonzero
endpoint coefficients choose both orthogonal types and a nonsingular
odd-dimensional radical. In the notation of that proof, an adjacent
product has characteristic polynomial
\[
                         X^2+\alpha^2X+1,\qquad\alpha\ne0.
\]
It has distinct reciprocal roots in $\mathbb F_q$ or
$\mathbb F_{q^2}$, and its order divides $q-1$ or $q+1$.
Thus every adjacent order is odd, including both free endpoints.
Contraction consequently gives all the rank intervals in survey
Theorem~10.3. The earlier stronger assertion that every chosen bond
has order exactly $q+1$ is unnecessary and is not retained for the
corrected endpoint construction.

\section{Suzuki groups, Ree groups and arbitrary diagrams}
The acquired Suzuki argument gives a particularly short example of
how subgroup structure and orbit counting complement one another.
The outer quotient of $\Aut(\Sz(2^N))$, for odd $N$, has odd order.
An involution-generated intermediate group therefore lies in the
simple socle. The acquired involution-centralizer and dihedral
arguments give precisely string rank three. Normalizing two commuting
endpoint involutions reduces the possible third marks to a field
parameter action. Exact subfield support, followed by field
Moebius inversion, removes the proper subfield groups. Orbit counting
and duality give precisely the formula in survey Theorem~6.2.
All Suzuki subgroup and generation theorems used in that derivation
remain explicit acquired inputs.

For small Ree groups, the complete acquired proof gives string rank
three for $R(3^N)$ with odd $N>1$, excludes larger ranks and excludes
the nonsimple $R(3)$. The odd outer quotient again excludes proper
almost simple extensions. The recovery replaces the source's false
normality step and keeps the exact subgroup and fusion hypotheses;
it does not transport the Suzuki proof by analogy.

Removing the string condition changes the problem. We retain as
published inputs the full arbitrary-C-group classification of Suzuki
groups by Connor--Leemans \cite{suzukiC}, including existence of a
nonstring marking, and the full $\PSL_2/\PGL_2$ classification by
Connor--Jambor--Leemans \cite{linearC}. The latter states that the
\emph{maximum} C-rank of $\PSL_2(q)$, $q\ge4$, is four exactly for
\[
                              q\in\{7,9,11,19,31\},
\]
and otherwise three; for $\PGL_2(q)$ it is four exactly at $q=5$.
\note{Correction}
Survey Theorem~11.2 instead says that a representation has rank four
if and only if its parameter is in this list. That literal universal
reading is false: $\PSL_2(11)$ and $\PGL_2(5)$ also have rank-three
string markings, hence rank-three C-markings. The cited authors'
publisher and conference abstracts explicitly state the maximal-rank
interpretation used here.

The two final symmetric C-group classifications are acquired in full.
For degree $n\ge7$ and rank $n-1$, the marks are precisely the edges
of a tree, so they are transpositions and generate $S_n$; the Coxeter
diagram is the line graph of that tree. For rank $n-2$ in $S_n$,
$n\ge9$, the three families of survey Theorem~11.4 are exactly those
of the acquired complete high-rank extension. Its necessity and
sufficiency arguments allow arbitrary branching in the attached tree.
The earlier incomplete version of that application is not used.

\section{The half-degree theorem and numerical conventions}
The acquired transitive-group theorem recovers all alternatives of
survey Theorem~10.1, including its complete degree-six and degree-eight
imprimitive exceptions and the primitive degree-six simplex. In the
infinite equality family, put $n=2m$ with $m$ odd. Write $\lambda$
for the simultaneous flip of all $m$ pairs, $e_1$ for one pair flip,
and $s_i$ for the adjacent transpositions of pairs. The actual tuple
\[
                     (\lambda,\lambda+e_1,s_1,\ldots,s_{m-1})
\]
has type $[2,4,3,\ldots,3]$. Reversing only the noncentral path gives
$[2,3,\ldots,3,4]$. These are generally distinct ordered polytopes
up to global duality, although their unordered distinguished set is
the same. The survey prints the latter type beside the former tuple.
Its Table~6 already has the correct $[3,2,3]$ type for $S_3\times S_3$;
we do not attribute the different typo in the earlier paper to this
survey edition.

For the simple groups and connected symmetric markings in Tables~1,
3 and~5, the noncommutation graph determines the path ordering up to
reversal. If a complete census has $c_r$ ordered classes and $s_r$
classes fixed by reversal, the survey convention is
\[
                            d_r=(c_r+s_r)/2.                  \tag{6}
\]
This follows by partitioning the classes into fixed points and
unordered two-element reversal orbits. It is not valid to halve
$c_r$ without also counting the self-dual classes.

The numerical appendix transcribes every numeric cell directly from
the fixed source TeX and records a disposition for each cell. The
completed L071 censuses recover every survey row through the Held
group, the entire $\mathrm{Co}_3$ row, the $S_7,S_9,A_5,A_9$ rows,
and the $A_{11},\ldots,A_{14}$ rows after applying (6). The exact
recorded self-dual counts are used, with no repeated search. The
symmetric stability theorem supplies every cell in its range, and
the complete maximal-rank theorems supply the zero cells above their
bounds. All fourteen rows of Table~2 follow from (5) and the
exceptional automorphism action. These are statements about named
cells, not a blanket validation of a table.

The complete retained $A_{10}$ lists have $177,4,8$ ordered classes
in ranks $3,4,5$. Direct reversal matching gives $11,0,0$ fixed
classes, hence the printed values $94,2,4$. The same calculation
on the complete $S_5$ list gives seven ordered rank-three classes
and one fixed class, hence four after reversal. The degree-six automorphism calculation gives the two rank-three
$S_6$ classes. The separately acquired complete six-point component
now also verifies its middle value, four in rank four, using the full
outer-automorphism action. The complete $S_8$ component gives ordered
counts $126,71,20,2,1$ and reversal fixed counts $10,1,2,0,1$.
Thus its source row is $68,36,11,1,1$; all four initial classes and
all $71\,296$ actual subset intersections are independently checked.
The retained $A_{10}$ and $S_5$ calculations reuse their completed
lists; the new $S_6$ and $S_8$ components execute the acquired common
method on precisely their fixed rows. The complete $S_{10}$ component
now gives ordered counts $809,403,99,26,13,2,1$ and reversal fixed counts
$17,3,5,0,1,0,1$. Its source row is therefore $413,203,52,13,7,1,1$.
All five initial involution classes finish, and a separate actual-group
verification checks every one of the $969\,024$ subset intersections.
The four entries in ranks three through six are newly acquired here;
the three higher entries were already consequences of stability.
Likewise, the complete $S_{11}$ calculation has ordered counts
$2382,377,86,49,18,13,2,1$ and reversal fixed counts
$60,1,0,1,0,1,0,1$. Its row is
$1221,189,43,25,9,7,1,1$. The complete five-class search and all
$3\,257\,472$ actual subset intersections have passed distinct review.
Only ranks three through six add new cells; the four higher entries
were already supplied by stability.

The last four rows are supplied by the companion \emph{Four finite rows
of the classification survey}. Its complete search has twenty-two initial
involution classes. Every branch terminates; its $94\,354$ actual tuples
are checked by $150\,173$ subgroup intersection tests. These tests use the
acquired interval criterion: a string marking has the full intersection
property if its two end subgroups do, and their intersection is exactly
the common middle subgroup. Induction over contiguous intervals, starting
with ranks one and two, proves the full property. Thus the reduction in
verification work is a theorem, not the omission of selected subset tests.
An independent literal conjugacy calculation pays the reversal convention.
The finite results are as follows; entries in each row start at rank three.
\begin{center}
\begin{tabular}{@{}cll@{}}
\toprule Group & Ordered classes & Classes modulo duality\\\midrule
$S_{12}$ & $6567,1875,364,150,74$ & $3346,940,183,75,40$\\
$S_{13}$ & $14181,1723,341,244,82$ & $7163,863,171,123,41$\\
$S_{14}$ & $45891,7869,1951,606,324,108$ & $23126,3945,978,303,163,54$\\
$A_{15}$ & $11557,314,40,81,12$ & $5820,158,20,42,6$\\\bottomrule
\end{tabular}
\end{center}
The reversal fixed counts are respectively
$(125,5,2,0,6)$, $(145,3,1,2,0)$, $(361,21,5,0,2,0)$ and
$(83,2,0,3,0)$. Higher symmetric ranks follow from (2), in its actual
stable range; the alternating maximum (3) excludes all ranks above seven.
Every one of the twenty-one previously unverified experimental cells is
therefore established without taking a target number as a premise.

There are $221$ displayed rank-count cells across Tables~1,~3 and~5.
Of these, $212$ now have an acquired project proof or complete census;
the remaining nine use the separately attributed published classification
in the next paragraph. This total does not subsume the different data in
the other tables. Table~2 has fourteen formula rows, Table~4 has the two
infinite constructions (4), and Table~6 has four finite group/type rows.
For the degree-six group of order $48$, Table~6 does not exhaust the
possible ordered types: the acquired list also gives $[2,4,3]$ and its
reversal. Each printed row is matched to a literal accepted tuple.

The Rudvalis, sporadic Suzuki and O'Nan entries refer to named
classification papers. The published Leemans--Mulpas paper~\cite{sporadic2022},
Section~5, states the complete counts $21594,227$ for Rudvalis,
$7119,257,13$ for Suzuki, and $6536,16$ for O'Nan, in successive ranks
starting at three, with no higher ranks. These are exact attributed
classification inputs; our separate recovery of that manuscript is not
yet complete. In particular, the literal ``Unknown'' in the survey's
O'Nan rank-three cell was answered in that 2022 publication. It is
preserved as historical source text, not presented as a presently open
question or as a count computed in this companion. The regular-map
count is a different quantity: the project's bounded correction to it
does not prove the O'Nan string count. Finally $|M_{22}|=443520$; Table~5's $443510$ is a
literal order typo, verified against the acquired exact group model.

\section{Questions, attribution and the recovery boundary}\label{questions}
Every one of the eighteen numbered problems or conjectures has an
individual record. Conjecture~1 is answered by the acquired later
stability theorem, with its original authorship retained. General
symmetric enumeration (Problem~1) is therefore partly answered, not
entirely solved. The construction and maximum-rank theorems for
alternating groups do not enumerate them (Problem~2).

The project already has internally reviewed written work on the
odd-$q$ maximum-four conjecture for $\PSL_4(q)$ (Conjecture~2), and
on part of Problem~7. Those exact earlier arguments, checks, reviews
and priority qualifications are linked as historical project results.
They are not assigned another conjecture credit here. Problem~7
includes field and graph extensions, whereas the recorded diagonal
interval results do not cover those extensions. The unitary companion
results are likewise not silently substituted for a linear-socle
question. No claim of historical priority follows from the earlier
bounded literature searches.

Problems~3--6 concern complete extension spectra or enumeration and
are not settled by the linear-socle obstruction. Problem~8 is a Ree
enumeration problem. Problem~9 asks for maximal ranks of the
characteristic-two classical groups; construction of the displayed
rank intervals alone does not prove maximality. Problem~10, the O'Nan rank-three computation, has the published answer
$6536$ in the 2022 classification just cited; it is not a new result of
this project. Problem~11 asks for the wider remaining sporadic classification
and is not settled by those three named groups alone. Problems~12--15 concern
arbitrary C-groups, whose questions are not answered by string
classifications. Conjecture~3, that $A_{11}$ is the only finite simple
group with a gap in its rank spectrum, is not proved by certifying
that one gap. These are dispositions of the fixed source questions,
not a new search asserting their present-day openness.

The new arguments in this survey integration include the companion's
complete cyclic-normal intersection proof, the explicit counterexamples
to the dimension-three and hypermap assertions, and the complete finite
classifications just stated.
Contraction, orbit multiplicities, pure-cut stability, supported
symmetry groups and the completed censuses are reused with their
previous authorship and reviews. The published external classification
premises are listed separately. Neither citations nor finite
agreement are treated as reproofs. Of the thirty-two numbered theorem and corollary clauses, twenty-five
have an acquired proof, correction or counterexample route, with their
stated prior inputs. The other seven---Theorems~3.1, 4.6, 5.1, 5.4, 5.6,
11.1 and 11.2---retain a published classification or exact construction.
Their statements and attribution are part of this synthesis; their proof
cost is not silently transferred to a newly proved theorem. The nine
sporadic numerical cells have the same explicit status. No additional
Lean claim or numerical compression factor is made at this intake stage.

\section*{Complete theorem index}
\small
\begin{longtable}{@{}>{\raggedright\arraybackslash}p{17mm}>{\raggedright\arraybackslash}p{44mm}>{\raggedright\arraybackslash}p{89mm}@{}}
\toprule Source & Acquired route or retained input & Exact qualification\\\midrule\endhead
3.1 & Simple-group obstruction; Nuzhin--Mazurov and corrections & Positive classification retained; all-rank negative obstruction acquired.\\
3.2 & Companion intersection theorem & Full no-cyclic-normal-subgroup hypothesis, including the Klein-four case.\\
4.1,4.3 & Contraction and path construction & All ranks; endpoint type corrected in (1).\\
4.2 & Pure cuts and 18 family models & Figure 1, stable range and connected diagram; exact source arrows matched.\\
4.4,4.10 & Orbit multiplicities & Exact formulas, parity and full degree-six automorphism fusion.\\
4.5 & Positive classification; alternating constructions & Rank-three existence, with retained published classification where used.\\
4.6 & Alternating interval; 2012 seeds & Existential conclusion acquired; the two exact types in (4) retained explicitly.\\
4.7--4.9 & Alternating interval and maximum theorem & Exact small groups and every rank in the stated interval; nontrivial degree domain.\\
5.1,5.4 & Sjerve--Cherkassoff & Complete rank-three classification retained; no substitution of a finite sample.\\
5.2,5.3,5.5 & Projective subgroup separators & Full rank-four classification with standard cell identifications retained.\\
5.6 & Connor--De Saedeleer--Leemans & Complete almost-simple classification is a named external input.\\
5.7 & Companion and dimension-three obstruction & Literal statement refuted; correct irreducible form obstruction retained.\\
5.8 & Simple-group obstruction; Brooksbank--Leemans & Even fields excluded; published odd-field rank-four construction retained.\\
6.1,6.2 & Suzuki normalization and field support & Exact rank theorem and both levels of Moebius inversion.\\
7.1 & Complete small-Ree proof & Actual subgroup/fusion hypotheses and repaired normality argument.\\
8.1 & Supported symmetries on quadratic paths & Both types and odd-dimensional quotient; corrected free endpoints.\\
10.1 & Complete transitive half-degree theorem & Literal wreath tuple/type distinguished from reordering and global duality.\\
10.2 & Actual subgroup contraction & Membership hypothesis retained.\\
Cor.1 & Exact contraction iteration & Tested bonds start with $p_3$; zero-step case always permitted.\\
10.3 & Quadratic paths plus contraction & All adjacent orders odd, sufficient for every stated rank.\\
11.1 & Connor--Leemans & Full arbitrary-C-group Suzuki theorem retained, including nonstring existence.\\
11.2 & Connor--Jambor--Leemans & Maximal C-rank classification retained; universal wording corrected.\\
11.3 & Extremal independent markings & Full tree necessity and sufficiency for $n\ge7$.\\
11.4 & Complete high-rank symmetric extension & All three families with arbitrary attached-tree branching.\\\bottomrule
\end{longtable}
\normalsize

\begin{thebibliography}{99}
\bibitem{survey} D. Leemans, \emph{String C-group representations of almost
simple groups: a survey}, arXiv:1910.08843v2, 11 May 2020; Contemp. Math. 764
(2021), 157--178. All source numbers in this paper refer to that fixed version.
\bibitem{stability} P. J. Cameron, M. E. Fernandes and D. Leemans,
\emph{The number of string C-groups of high rank}, fixed
arXiv:2212.12723v2. Acquired complete corrected principal recovery L019.
\bibitem{alt2012} M. E. Fernandes, D. Leemans and M. Mixer,
\emph{Polytopes of high rank for the alternating groups},
J. Combin. Theory Ser. A 119 (2012), 42--56, Theorem 1.1. The exact seeds in (4)
remain inputs, including generation and full intersection property.
\bibitem{SC} D. Sjerve and M. Cherkassoff, \emph{On groups generated by
three involutions, two of which commute}, CRM Proc. Lecture Notes 6
(1994), 169--185. Retained precise existence classification.
\bibitem{almost2} T. Connor, J. De Saedeleer and D. Leemans,
\emph{Almost simple groups with socle PSL(2,q) acting on abstract regular
polytopes}, J. Algebra 423 (2015), 550--558. Retained full classification.
\bibitem{linear4} P. A. Brooksbank and D. Leemans,
\emph{Polytopes of large rank for PSL(4,q)}, J. Algebra 452 (2016), 390--400,
Theorem 1.1. Retained odd-field rank-four existence.
\bibitem{suzukiC} T. Connor and D. Leemans, \emph{C-groups of Suzuki type},
J. Algebraic Combin. 42 (2015), 849--860, Classification Theorem 1.1.
\bibitem{linearC} T. Connor, S. Jambor and D. Leemans,
\emph{C-groups of PSL(2,q) and PGL(2,q)}, J. Algebra 427 (2015), 455--466,
Theorem 1.1. Retained maximal-rank and explicit rank-four classification.
\bibitem{hypermaps} M. Conder, P. Poto\v cnik and J. \v Sir\'a\v n,
\emph{Regular hypermaps over projective linear groups},
J. Aust. Math. Soc. 85 (2008), 155--175. The inner-reflexible correspondence
is retained; the blanket transfer to C-group markings is refuted.
\bibitem{sporadic2022} D. Leemans and J. Mulpas,
\emph{The string C-group representations of the Suzuki, Rudvalis and
O'Nan sporadic groups}, Art Discrete Appl. Math. 5 (2022),\#P3.09,
Section~5; doi:10.26493/2590-9770.1405.4ce. Exact external census inputs;
no independent re-exhaustion is claimed in this survey companion.
\bibitem{censuses} D. Leemans and M. Mixer, \emph{Algorithms for
classifying regular polytopes with a fixed automorphism group},
Contrib. Discrete Math. 7 (2012), 105--118. The complete acquired L071
recovery includes the earlier Hartley--Hulpke lists and exact duality counts.
\end{thebibliography}
\end{document}
